\documentclass[11pt]{article}
\usepackage[a4paper,margin=25mm]{geometry}
\usepackage{amsmath,amssymb,amsthm,booktabs,array}
\usepackage[T1]{fontenc}
\usepackage{lmodern}
\usepackage[colorlinks=true,urlcolor=blue,citecolor=blue,linkcolor=blue]{hyperref}
\hypersetup{pdftitle={Alikhani--Griswold, Section 4(3): connected graphs with domination root -4},pdfauthor={Haobo Ma; Wenlin Zhang},pdfsubject={Connected domination-root families; CC BY 4.0 (cc-by-4.0)}}
\newtheorem{theorem}{Theorem}
\newtheorem{lemma}[theorem]{Lemma}
\newtheorem{corollary}[theorem]{Corollary}
\newcommand{\Dom}{\mathcal D}
\newcommand{\base}{H}
\title{Alikhani--Griswold, Section 4(3):\\connected graphs with domination root $-4$}
\author{Haobo Ma\thanks{ChronoAI Pte Ltd, Singapore.}\and
Wenlin Zhang\thanks{National University of Singapore; ChronoAI Pte Ltd, Singapore.}}
\date{2 November 2026}
\begin{document}
\maketitle
\begin{abstract}
We answer Section 4, item 3 of Alikhani and Griswold's
\emph{On the Integer Domination Root Conjecture} by constructing connected
graphs with domination root $-4$ and unbounded order.
Joining $r$ copies of their 33-vertex graph through vertices adjacent to
leaves gives a connected graph on $33r$ vertices whose domination polynomial
is the $r$th power of the original polynomial. The root $-4$ has multiplicity
exactly $r$, and the copies can be joined along a path with maximum degree
seven. A mixed construction gives examples of every order at least 35,
with maximum degree seven and cycle rank four.
We give a complete set-bijection proof, an exact finite calculation
for the base graph, and the scope of the corresponding Lean declarations.
\end{abstract}

\section{The question and the construction}
Alikhani and Griswold~\cite[Section 4, item 3]{AG} ask whether infinitely
many graphs have ordinary domination polynomial zero at $-4$.
Their convention is finite simple undirected connected graphs.
For any finite simple graph $G$, write
\[
 \Dom(G)=\{S\subseteq V(G):N_G[S]=V(G)\},\qquad
 D(G,x)=\sum_{S\in\Dom(G)}x^{|S|}.
\]
Connectedness is required of $G$, not of the sets $S$.
The source supplies a connected graph $\base$ of order 33 with
$D(\base,-4)=0$~\cite[Theorem 2.1 and Corollary 2.2]{AG}.
Its vertices 4 and 5 satisfy $N_\base(5)=\{4\}$.

The construction uses the established theory of irrelevant edges:
Kotek, Preen, Simon, Tittmann and Trinks characterize such edges using
domination-covered vertices~\cite[Definition 3.9 and Theorem 3.11]{KPSTT}.
A vertex with a leaf neighbor is domination-covered. Consequently,
adding an edge between two nonadjacent vertices that each have a leaf
neighbor preserves the domination polynomial.
The following direct proof records exactly the special case needed here.

\begin{lemma}\label{lem:join}
Let $H_1,\ldots,H_r$ be vertex-disjoint finite simple graphs, each with
specified vertices $v_i,w_i$ satisfying $N_{H_i}(w_i)=\{v_i\}$.
Let $F$ be any simple graph on $\{1,\ldots,r\}$.
Obtain $J_F$ from their disjoint union by adding $v_iv_j$ for each
$ij\in E(F)$. Then
\begin{equation}\label{eq:product}
 \Dom(J_F)=\{S_1\cup\cdots\cup S_r:S_i\in\Dom(H_i)\},\qquad
 D(J_F,x)=\prod_{i=1}^rD(H_i,x).
\end{equation}
If every $H_i$ and $F$ are connected, then $J_F$ is connected.
\end{lemma}
\begin{proof}
Let $S$ dominate $J_F$, and put $S_i=S\cap V(H_i)$.
The leaf $w_i$ has no neighbor outside $H_i$, so $v_i\in S_i$ or
$w_i\in S_i$. In either case $v_i$ is dominated within $H_i$.
Every vertex of $H_i$ other than $v_i$ has exactly its original
neighbors, so its domination by $S$ also occurs within $H_i$.
Thus $S_i$ dominates $H_i$.
Conversely, the union of internally dominating sets $S_i$ dominates
$J_F$, since none of the original edges was removed.
Restriction and union are inverse and preserve cardinality as
$|S|=\sum_i|S_i|$. Summing $x^{|S|}$ gives~\eqref{eq:product} as an
identity in $\mathbb Z[x]$.
For connectedness, every vertex of $H_i$ has a path to $v_i$, and
the added edges give paths between the $v_i$ according to $F$.
\end{proof}

\begin{theorem}\label{thm:family}
For every integer $r\geq1$ there is a connected finite simple graph
$G_r$ on $33r$ vertices with
\[
 D(G_r,x)=D(\base,x)^r.
\]
The root $-4$ has multiplicity exactly $r$. The graphs can be chosen
with $37r-1$ edges and maximum degree seven.
Consequently, for every $N\in\mathbb N$ there is such a graph of order
strictly greater than $N$.
\end{theorem}
\begin{proof}
Use $r$ disjoint copies of $\base$, take $v_i$ and $w_i$ to be the
copies of 4 and 5, and use the path $F=P_r$ in Lemma~\ref{lem:join}.
The polynomial identity and connectedness follow immediately.
The source's simple root at $-4$, also verified in Section~\ref{sec:base},
has multiplicity $r$ in its $r$th power.
The base has 36 edges, and the path adds $r-1$, giving $37r-1$.
Its maximum degree is seven, attained at vertex 12; only the copied
vertices 4 acquire new neighbors, and their degrees rise from two
to at most four. Thus the maximum degree remains seven.
Taking $r=N+1$ gives $33r>N$.
Distinct values of $r$ give distinct orders, hence nonisomorphic graphs.
\end{proof}

\section{An exact calculation for the base graph}\label{sec:base}
We use the factual edge list in~\cite[Section 2]{AG}.
For a self-contained specification, its edges are
\begin{align*}
E(\base)=\{&01,03,04,0\!:\!12,12,1\!:\!12,1\!:\!18,23,26,28,2\!:\!10,2\!:\!18,\\
&2\!:\!21,3\!:\!21,3\!:\!24,45,67,89,10\!:\!11,12\!:\!13,12\!:\!14,\\
&12\!:\!15,12\!:\!16,12\!:\!17,18\!:\!19,18\!:\!20,21\!:\!22,21\!:\!23,\\
&24\!:\!25,24\!:\!26,26\!:\!27,27\!:\!28,28\!:\!29,28\!:\!31,29\!:\!30,31\!:\!32\}.
\end{align*}
Here $ab$ for single-digit labels means $\{a,b\}$; a colon separates
endpoints when at least one label has two digits.
Remove the core $C=\{0,1,2,3\}$, whose induced graph is the cycle
$01,12,23,30$. The resulting eight components $B_i$ have roots,
vertex intervals and core-neighbor sets $I_i$ as in Table~\ref{tab:messages}.
Each root has a leaf neighbor in its own component.
For $b\in\{0,1\}$, define
\[
 M_i(b;x)=\sum_{\substack{T\in\Dom(B_i)\\\boldsymbol1_{a_i\in T}=b}}x^{|T|},
\]
where $a_i$ is the root. For $Z\subseteq C$ and
$A\subseteq\{1,\ldots,8\}$, let $c(Z,A)=1$ when every core vertex is
in $Z$, adjacent in $C$ to $Z$, or in $I_i$ for some $i\in A$;
otherwise let $c(Z,A)=0$. Then
\begin{equation}\label{eq:core}
 D(\base,x)=\sum_{Z\subseteq C}x^{|Z|}
       \sum_{A\subseteq\{1,\ldots,8\}}c(Z,A)
            \prod_{i=1}^8 M_i(\boldsymbol1_{i\in A};x).
\end{equation}
To justify the formula, restrict a dominating set to $C$ and each $B_i$.
All nonroot component vertices must be dominated internally because
their neighbors are internal. The component leaf forces its root to be
selected or adjacent to the selected leaf, so the root is internally
dominated too. The remaining condition is exactly $c(Z,A)=1$, with
$A$ recording selected roots. Conversely these conditions ensure
global domination. Restriction and disjoint union are inverse, and
the subset sizes add, proving~\eqref{eq:core}.

\begin{table}[ht]
\centering
\begin{tabular}{@{}rrrrr@{}}
\toprule
Root $a_i$ & Vertices of $B_i$ & $I_i$ & $M_i(0;-4)$ & $M_i(1;-4)$\\
\midrule
4 & $4,5$ & $\{0\}$ & $-4$ & $12$\\
6 & $6,7$ & $\{2\}$ & $-4$ & $12$\\
8 & $8,9$ & $\{2\}$ & $-4$ & $12$\\
10 & $10,11$ & $\{2\}$ & $-4$ & $12$\\
12 & $12,\ldots,17$ & $\{0,1\}$ & $-1024$ & $972$\\
18 & $18,19,20$ & $\{1,2\}$ & $16$ & $-36$\\
21 & $21,22,23$ & $\{2,3\}$ & $16$ & $-36$\\
24 & $24,\ldots,32$ & $\{3\}$ & $5888$ & $-20736$\\
\bottomrule
\end{tabular}
\caption{Root-conditioned domination sums for the source graph.}
\label{tab:messages}
\end{table}

Enumeration of the $608=4\cdot2^2+2^6+2\cdot2^3+2^9$ component
subsets gives Table~\ref{tab:messages}; the $2^4 2^8=4096$ core/root
states in~\eqref{eq:core} then give $D(\base,-4)=0$.
Grouping contributions by $|Z|=0,1,2,3,4$ gives, respectively,
\begin{gather*}
27135603376128,\quad -148575803670528,\quad 121440200294400,\\
-323840534118400,\quad 323840534118400,
\end{gather*}
whose sum is zero. Keeping the cardinality weights as polynomials
in~\eqref{eq:core} reproduces every coefficient in the source's
Table 1 and its factorization
\[
 D(\base,x)=x^{11}(x+2)(x+4)(x^2+3x+1)Q_{18}(x).
\]
In particular $D'(\base,-4)=1100417597440000\ne0$, so the root is simple.
The polynomial $Q_{18}$ is monic with constant coefficient 36.
Its evaluations at all 18 signed divisors of 36 are nonzero, so
the rational root theorem excludes integer roots of $Q_{18}$.
The quadratic factor has discriminant five. Thus the integer roots of $D(\base,x)$,
and hence of $G_r$ and $R_r$ for $r\geq1$, are exactly $0,-2,-4$.
The accompanying \texttt{graph.json} and standard-library Python program
\texttt{verify.py} implement this finite calculation with exact integers.
They also compare complete domination polynomials for 97 small support
joins, with at most 12 vertices, against~\eqref{eq:product}.

\section{Consequences and verification scope}
The same base yields a second explicit family $R_r$: add a new vertex
$h$ adjacent to vertex 4 in each of $r$ copies of $\base$, and add a
new leaf $p$ adjacent to $h$. For $r\geq1$,
\begin{equation}\label{eq:hub}
 |V(R_r)|=33r+2,\qquad D(R_r,x)=x(x+2)D(\base,x)^r.
\end{equation}
Indeed, Lemma~\ref{lem:join} applies to the $r$ copies and the extra
component $hp$, joining their support vertices 4 and $h$ along a star.
Since $D(K_2,x)=x(x+2)$, it gives~\eqref{eq:hub} and connectedness.
For $r=0$ the graph is $K_2$, with $D(K_2,-4)=8$.

\begin{corollary}\label{cor:orders}
Every integer $n\geq35$ is the order of a connected finite simple graph
with maximum degree seven, cycle rank four and simple domination root
$-4$. Its integer domination roots are exactly $0,-2,-4$.
\end{corollary}
\begin{proof}
Take $\base$, $a\geq0$ copies of $K_2$ and $b\geq0$ copies of $P_3$.
Use vertex 4 in $\base$, one endpoint in each $K_2$, and the middle
vertex in each $P_3$ as supports. Join these supports along a path.
Lemma~\ref{lem:join} gives a connected graph $T_{a,b}$ with
\begin{align*}
 |V(T_{a,b})|&=33+2a+3b, & |E(T_{a,b})|&=36+2a+3b,\\
 D(T_{a,b},x)&=D(\base,x)[x(x+2)]^a[x(x^2+3x+1)]^b.
\end{align*}
Here $D(P_3,x)=x+3x^2+x^3$, by listing its dominating subsets.
The added factors have values 8 and $-20$ at $-4$, respectively,
so that root remains simple; they introduce no new integer roots.
The cycle rank is $|E|-|V|+1=4$.
New supports have degree at most four, and the maximum degree seven
of $\base$ is unchanged.
For odd $n\geq35$ take $b=0$ and $a=(n-33)/2$; for even $n\geq36$
take $b=1$ and $a=(n-36)/2$.
\end{proof}

The pinned Lean file~\cite{Lean} contains three relevant declarations.
Their full namespace is
\begin{flushleft}\small
\nolinkurl{D5.S3.Combinatorics.Graph.Domination.MinusFourReassembly}.
\end{flushleft}
\begin{itemize}
\item \texttt{reassembly} proves the general weighted core/component
identity for every decomposition satisfying its stated neighbor and leaf
conditions and every integer weight.
\item \texttt{zero\_of\_messages} transfers exact component-message
equalities and a zero core sum to a zero domination evaluation.
\item \texttt{result} proves that for every $N\in\mathbb N$ there are
$n>N$ and a connected simple graph on $\mathrm{Fin}(n)$ whose ordinary
dominating-set sum at $-4$ is zero. Its witnesses use~\eqref{eq:hub}.
\end{itemize}
The polynomial identity~\eqref{eq:product}, the $33r$-vertex family,
its degree bound and multiplicity statement, and Corollary~\ref{cor:orders}
have the written proofs
above. Table~\ref{tab:messages}, the full base polynomial and the 97
small joins are exact finite checks. These scopes distinguish the
written polynomial results from the integer-evaluation Lean statements.
The declarations are linked directly in the
\href{https://github.com/the-omega-institute/trureturing/blob/a87bdf39e4517497a9773c0c724243bc700c5bb1/D5/S3/Combinatorics/Graph/Domination/MinusFourReassembly.lean}{pinned \texttt{MinusFourReassembly.lean} source}.

\begin{center}
\begin{tabular}{@{}p{.33\linewidth}p{.60\linewidth}@{}}
\toprule
Source item & Coverage\\
\midrule
Section 4(3), infinite family & Proved by Theorem~\ref{thm:family}; the
Lean \texttt{result} proves the same unbounded-order conclusion.\\[3pt]
Section 4(1), other integer roots & Not settled in general.
For $r\geq1$, $G_r$ and $R_r$ have exactly the integer roots $0,-2,-4$,
as does $T_{a,b}$ for $a,b\geq0$.\\[3pt]
Section 4(2), minimum order & Not covered by this argument; no
minimum-order assertion is used.\\[3pt]
Corollary 2.2, simple base root & Prior result of Alikhani--Griswold,
confirmed by the exact base derivative above.\\
\bottomrule
\end{tabular}
\end{center}

\paragraph{Data and code.}
The source package contains the LaTeX note, graph data, finite checker,
its output, metadata and build instructions.
The base graph is attributed to~\cite{AG}; the
finite checker is an independent implementation of~\eqref{eq:core}.

\begin{thebibliography}{9}
\raggedright
\bibitem{AG} S. Alikhani and M. Griswold,
\emph{On the Integer Domination Root Conjecture},
arXiv:2608.00109v1 (2026).
\href{https://arxiv.org/abs/2608.00109v1}{arXiv:2608.00109v1}.
\bibitem{KPSTT} T. Kotek, J. Preen, F. Simon, P. Tittmann and M. Trinks,
\emph{Recurrence relations and splitting formulas for the domination polynomial},
Electronic Journal of Combinatorics \textbf{19}(3) (2012), P47.
\href{https://doi.org/10.37236/2475}{doi:10.37236/2475}.
\bibitem{Lean} The Omega Institute,
\emph{MinusFourReassembly.lean}, revision
\texttt{a87bdf39e4517497a9773c0c724243bc700c5bb1}.
\href{https://github.com/the-omega-institute/trureturing/blob/a87bdf39e4517497a9773c0c724243bc700c5bb1/D5/S3/Combinatorics/Graph/Domination/MinusFourReassembly.lean}{Exact source and declarations}.
\end{thebibliography}
\noindent\small This note is licensed under Creative Commons Attribution 4.0
International (CC BY 4.0, \texttt{cc-by-4.0}).
\end{document}
